\section{4. Especificación de requisitos de software
(SRS)}\label{especificaciuxf3n-de-requisitos-de-software-srs}

Estructura adaptada de ISO/IEC/IEEE 29148:2018. Cada requisito tiene un
identificador único que se utiliza para la trazabilidad.

\subsection{4.1 Propósito, alcance y audiencia del
SRS}\label{propuxf3sito-alcance-y-audiencia-del-srs}

\textbf{Propósito:} este SRS describe qué debe hacer el Sistema de
Órdenes de Servicio y Citas y con qué calidad, para que el cliente lo
valide y lo firme en la Review 2, y para que el equipo lo use como base
del diseño, la implementación y las pruebas.

\textbf{Alcance:} cubre el ciclo de atención del taller (citas,
recepción, diagnóstico, cotización y autorización, refacciones, cobro y
entrega) y el portal web del cliente, según el alcance de la sección
3.3.

\textbf{Audiencia:} los dueños y el gerente del taller (aprobación), el
personal que usará el sistema (validación de sus flujos), el equipo de
desarrollo y el docente evaluador.

\textbf{Convenciones:} cada requisito tiene un identificador único: RF
(funcional), RNF (no funcional), RN (regla de negocio), HU (historia de
usuario) y CU (caso de uso). Estos identificadores se usan en la matriz
de trazabilidad (5.3).

\subsection{4.2 Glosario del dominio}\label{glosario-del-dominio}

{\def\LTcaptype{none} % do not increment counter
\begin{longtable}[]{@{}
  >{\raggedright\arraybackslash}p{(\linewidth - 4\tabcolsep) * \real{0.2349}}
  >{\raggedright\arraybackslash}p{(\linewidth - 4\tabcolsep) * \real{0.5342}}
  >{\raggedright\arraybackslash}p{(\linewidth - 4\tabcolsep) * \real{0.2309}}@{}}
\toprule\noalign{}
\begin{minipage}[b]{\linewidth}\raggedright
\textbf{Término}
\end{minipage} & \begin{minipage}[b]{\linewidth}\raggedright
\textbf{Definición}
\end{minipage} & \begin{minipage}[b]{\linewidth}\raggedright
\textbf{Sinónimos que usa el cliente}
\end{minipage} \\
\midrule\noalign{}
\endhead
\bottomrule\noalign{}
\endlastfoot
Autorización & Aceptación de una cotización por el cliente; registra
fecha, hora, IP, monto y el código de verificación. & Aprobación, firma
digital \\
Bahía & Espacio de trabajo donde se repara un vehículo. El taller tiene
4 y cada una aloja un vehículo a la vez. & --- \\
Checklist de inspección & Lista de verificación del estado de
carrocería, interiores y accesorios al recibir el vehículo. & Lista de
inspección \\
Cita & Solicitud para atender un vehículo en una fecha y hora; queda
«Pendiente» hasta que la recepcionista la confirma. & --- \\
Código 2FA & Código de 6 dígitos, de un solo uso y vigente 5 minutos,
que se envía al correo del cliente para confirmar una autorización. &
Código de verificación, doble verificación \\
Compra & Adquisición de refacciones a un proveedor, respaldada por un
ticket. & Ticket de compra \\
Concepto & Renglón de una cotización: mano de obra, refacción o insumo.
Puede ser opcional. & Rubro \\
Cotización & Costo desglosado de mano de obra, refacciones e insumos de
una orden; vigente 5 días naturales. & Presupuesto \\
Diagnóstico & Descripción técnica de la falla y de las reparaciones
necesarias que registra el jefe de mecánicos. & --- \\
Enlace seguro & Liga web única que recibe el cliente para consultar su
vehículo y autorizar cotizaciones sin usuario ni contraseña. & Token de
acceso \\
Entrega a crédito & Entrega del vehículo con saldo pendiente; solo para
flotillas registradas y con autorización del gerente. & Crédito \\
Flotilla & Empresa cliente con varios vehículos que puede recibir
entrega a crédito. & Flotilla empresarial \\
Folio & Identificador único de una orden de servicio. & --- \\
Fotografía de recepción & Evidencia digital del estado del vehículo al
llegar. Son obligatorias cuatro: frente, trasera, costado izquierdo y
costado derecho. & --- \\
Garantía & Cobertura de 30 días sobre la mano de obra, contados desde la
entrega del vehículo. & --- \\
Insumo & Material consumible que se usa en la reparación, como aceite o
líquidos. & Consumible \\
Jefe de mecánicos & Líder técnico: diagnostica, cotiza, asigna bahías y
mecánicos y solicita refacciones desde la tableta. & Líder técnico \\
Mecánico responsable & Mecánico titular de una orden; otros mecánicos
pueden colaborar sin cambiar la titularidad. & Mecánico principal \\
Orden de servicio & Registro de todo el trabajo sobre un vehículo, desde
la recepción hasta la entrega. & Orden \\
Proveedor & Refaccionaria local a la que el almacén compra piezas. &
Refaccionaria \\
Recepción vehicular & Registro de la llegada del vehículo: kilometraje,
nivel de combustible, checklist y fotografías. & Check-in, registro de
entrada \\
Re-cotización & Nueva versión de la cotización cuando los trabajos
adicionales aumentan el monto más de 10 \%. & --- \\
Refacción & Pieza que se instala en el vehículo. & Pieza, repuesto \\
Saldo pendiente & Diferencia entre el total autorizado de la orden y lo
que el cliente ha pagado. & --- \\
Solicitud de refacción & Pedido de piezas que el jefe de mecánicos hace
al almacén para una orden. & Requisición \\
\end{longtable}
}

\subsection{4.3 Descripción general}\label{descripciuxf3n-general}

\subsubsection{4.3.1 Usuarios y sus
características}\label{usuarios-y-sus-caracteruxedsticas}

El sistema será utilizado por distintos perfiles dentro de la estructura
operativa del negocio y la parte gerencial del taller mecánico, además
de los clientes finales del negocio. La siguiente tabla detalla las
características, nivel técnico y contexto de uso:

{\def\LTcaptype{none} % do not increment counter
\begin{longtable}[]{@{}
  >{\raggedright\arraybackslash}p{(\linewidth - 6\tabcolsep) * \real{0.2138}}
  >{\raggedright\arraybackslash}p{(\linewidth - 6\tabcolsep) * \real{0.2990}}
  >{\raggedright\arraybackslash}p{(\linewidth - 6\tabcolsep) * \real{0.2138}}
  >{\raggedright\arraybackslash}p{(\linewidth - 6\tabcolsep) * \real{0.2734}}@{}}
\toprule\noalign{}
\begin{minipage}[b]{\linewidth}\raggedright
\textbf{Tipo de usuario}
\end{minipage} & \begin{minipage}[b]{\linewidth}\raggedright
\textbf{Tareas principales}
\end{minipage} & \begin{minipage}[b]{\linewidth}\raggedright
\textbf{Experiencia tecnológica}
\end{minipage} & \begin{minipage}[b]{\linewidth}\raggedright
\textbf{Dispositivo y contexto de uso}
\end{minipage} \\
\midrule\noalign{}
\endhead
\bottomrule\noalign{}
\endlastfoot
Recepcionista & Registro de clientes/vehículos, check-in con fotos,
apertura de órdenes, confirmación de citas web, cobranza y entrega. &
Media: Manejo fluido de navegación web, correo electrónico y paquetes de
oficina. & Computadora de escritorio en mostrador de recepción. Entorno
de oficina con teclado/mouse. \\
Jefe de Mecánicos (Líder Técnico) & Diagnóstico técnico, elaboración de
cotizaciones, asignación de bahías/mecánicos, requisición de refacciones
y sellado final. & Baja--Media: Familiarizado con smartphones
personales; poca experiencia con sistemas de gestión. & Tabletas
táctiles exclusivas de taller (1 o 2 dispositivos). Entorno con ruido,
polvo y manipulación con manos sucias/grasosas. \\
Mecánico en Bahía & Ejecución de reparaciones físicas, notificación de
hallazgos adicionales o solicitud de refacciones al Jefe de Mecánicos. &
Baja: Uso de teléfonos inteligentes para uso personal cotidiano. & Sin
dispositivo asignado directo; interactúa a través del Jefe de Mecánicos
para evitar deterioro de hardware \\
Encargado de Almacén & Atender solicitudes de piezas, consultar
directorio de proveedores locales, registrar compras/tickets e historial
de insumos. & Baja--Media\textbf{:} Uso básico de computadora para
consulta de inventarios. & Computadora de escritorio o terminal dedicada
en el área de almacén/refacciones. \\
Gerente (Dueño) & Monitoreo del Dashboard de KPIs (ocupación, ingresos),
aprobación de entregas a crédito a flotillas y excepciones. & Media:
Manejo de smartphones y computadoras para administración del negocio. &
Laptop personal o dispositivo móvil (smartphone/tableta) con acceso
remoto/local \\
Cliente & Solicitud de citas en línea, seguimiento en vivo del servicio
y aprobación/rechazo de cotizaciones mediante código 2FA. & Variable
(Baja a Alta): Desde usuarios sin conocimientos técnicos hasta usuarios
avanzados. & Dispositivo móvil personal (iOS/Android) vía navegador web
progresivo (PWA) sin instalar aplicaciones. \\
\end{longtable}
}

\subsubsection{4.3.2 Restricciones, supuestos y
dependencias}\label{restricciones-supuestos-y-dependencias}

\begin{enumerate}
\def\labelenumi{\arabic{enumi}.}
\item
  \textbf{Restricciones:}
\end{enumerate}

\begin{itemize}
\item
  Presupuesto y Licenciamiento: El desarrollo debe realizarse
  estrictamente con herramientas, lenguajes y frameworks de código
  abierto (\emph{Open Source} como Python y FastAPI), garantizando cero
  costos de licencias o rentas de software.
\item
  Infraestructura Física: El taller cuenta únicamente con 4 bahías de
  trabajo físicas (capacidad estricta de 1 vehículo por bahía) y un
  estacionamiento con capacidad máxima de 5 vehículos en espera.
\item
  Dispositivos Hardware: La captura digital en taller se limita a máximo
  2 tabletas táctiles centralizadas por el Jefe de Mecánicos para
  proteger la inversión y evitar daños.
\end{itemize}

\begin{itemize}
\item
  Alcance del Vehículo\textbf{:} La plataforma atenderá exclusivamente
  vehículos ligeros (automóviles particulares, motocicletas y
  camionetas), excluyendo maquinaria pesada o industrial.
\end{itemize}

\begin{enumerate}
\def\labelenumi{\arabic{enumi}.}
\setcounter{enumi}{1}
\item
  \textbf{Supuestos:}
\end{enumerate}

\begin{itemize}
\item
  Se asume que el taller mantiene activa una conexión de red local
  (Wi-Fi) para la sincronización entre computadoras de mostrador y
  tabletas de taller en condiciones normales.
\item
  Se asume que los clientes cuentan con un correo electrónico para
  recibir el enlace de seguimiento y los códigos de verificación de 6
  dígitos (2FA); el número celular se usa para el envío opcional del
  enlace por SMS (RF-02.3).
\item
  Se asume que los datos ingresados durante la fase de desarrollo y
  pruebas son 100\% ficticios para dar cumplimiento a la Ley Federal de
  Protección de Datos Personales.
\end{itemize}

\begin{enumerate}
\def\labelenumi{\arabic{enumi}.}
\setcounter{enumi}{2}
\item
  \textbf{Dependencias:}
\end{enumerate}

\begin{itemize}
\item
  Infraestructura Eléctrica y Red: Operación dependiente de baterías en
  laptops/tabletas y arquitectura \emph{Local-First / PWA} (Service
  Workers e IndexedDB) para permitir la captura de datos \emph{offline}
  durante cortes de luz o fallas de internet.
\item
  Proveedores Externos de Comunicación: Dependencia del servicio de
  correo electrónico para entregar enlaces y códigos 2FA al cliente, y
  de un servicio de SMS para el envío opcional del enlace.
\item
  Disponibilidad de Refaccionarias Locales: La adquisición de piezas
  fuera de stock depende del directorio de proveedores locales e
  historial de compras del almacén.
\end{itemize}

\subsection{4.4 Requisitos funcionales}\label{requisitos-funcionales}

{\def\LTcaptype{none} % do not increment counter
\begin{longtable}[]{@{}
  >{\raggedright\arraybackslash}p{(\linewidth - 8\tabcolsep) * \real{0.0962}}
  >{\raggedright\arraybackslash}p{(\linewidth - 8\tabcolsep) * \real{0.4444}}
  >{\raggedright\arraybackslash}p{(\linewidth - 8\tabcolsep) * \real{0.1174}}
  >{\raggedright\arraybackslash}p{(\linewidth - 8\tabcolsep) * \real{0.1497}}
  >{\raggedright\arraybackslash}p{(\linewidth - 8\tabcolsep) * \real{0.1923}}@{}}
\toprule\noalign{}
\begin{minipage}[b]{\linewidth}\raggedright
\textbf{ID}
\end{minipage} & \begin{minipage}[b]{\linewidth}\raggedright
\textbf{Requisito funcional}
\end{minipage} & \begin{minipage}[b]{\linewidth}\raggedright
\textbf{Módulo}
\end{minipage} & \begin{minipage}[b]{\linewidth}\raggedright
\textbf{Prioridad (MoSCoW)}
\end{minipage} & \begin{minipage}[b]{\linewidth}\raggedright
\textbf{Historias relacionadas}
\end{minipage} \\
\midrule\noalign{}
\endhead
\bottomrule\noalign{}
\endlastfoot
RF-01 & El sistema deberá permitir autenticar mediante credenciales
seguras al personal del taller según su rol (Gerente, Recepcionista,
Jefe de Mecánicos y Encargado de Almacén). & M1 -- Autenticación y
Seguridad & Must & HU-27 \\
RF-02.1 & El sistema deberá permitir al cliente consultar el estado de
su vehículo y autorizar cotizaciones mediante un enlace web seguro
enviado por correo electrónico o por SMS (RF-02.3), sin instalar
aplicaciones. & M1 -- Autenticación y Seguridad & Must & HU-02, HU-28 \\
RF-02.2 & El sistema deberá de expirar el enlace del cliente de manera
automática una vez que el vehículo sea entregado y saldado. & M1 --
Autenticación y Seguridad & Should & HU-02, HU-23 \\
RF-02.3 & El sistema deberá de permitir enviar el enlace del cliente
también por SMS. & M1 -- Autenticación y Seguridad & Could & HU-02 \\
RF-03.1 & El sistema deberá registrar, buscar y actualizar la
información de clientes y vehículos (placas, marca, modelo, año y VIN
(opcional)). & M3 -- Recepción E Historial Inicial & Must & HU-05,
HU-06 \\
RF-03.2 & El sistema deberá mantener el historial de servicios de cada
vehículo, con fotografías de recepción, diagnósticos técnicos y
garantías aplicables (30 días en mano de obra). & M7 -- Portal del
Cliente, Cobro y Gerencia & Should & HU-23, HU-24 \\
RF-04 & El sistema deberá registrar en la recepción el kilometraje, el
nivel de combustible y un checklist de carrocería, y exigir fotografías
digitales del estado inicial de la unidad. & M3 -- Recepción E Historial
Inicial & Must & HU-01, HU-25, HU-26 \\
RF-05 & El sistema deberá permitir abrir órdenes de servicio asociadas a
una de las 4 bahías y a un mecánico responsable, controlando los estados
operativos del taller. & M5 -- Taller, Bahías y Expediente Técnico &
Must & HU-07, HU-08, HU-09, HU-11 \\
RF-06 & El sistema deberá generar cotizaciones desglosadas (mano de obra
y refacciones/insumos) con vigencia de 5 días naturales y requerir firma
digital validada con un código enviado al correo del cliente, de forma
remota o presencial. & M4 Cotización y Contratos 2FA & Must & HU-02,
HU-12, HU-13, HU-14, HU-15 \\
RF-07 & El sistema deberá bloquear el avance del trabajo y exigir una
nueva autorización con verificación cuando los trabajos adicionales
incrementen el monto inicial en más de 10 \%. & M4 Cotización y
Contratos 2FA & Must & HU-16 \\
RF-08.1 & El sistema deberá permitir al jefe de mecánicos solicitar al
almacén las piezas que requieran los mecánicos, vinculadas a una orden
de servicio. & M6 -- Almacén y Refacciones & Should & HU-10, HU-19 \\
RF-08.2 & El sistema deberá permitir al encargado de almacén mantener un
directorio de proveedores locales (teléfonos o correos) y registrar
costos y fechas de llegada de las piezas. & M6 -- Almacén y Refacciones
& Should & HU-17, HU-18 \\
RF-08.3 & El sistema deberá de emitir un ticket e historial de compras
asociado a cada orden de servicio. & M6 -- Almacén y Refacciones & Could
& HU-18 \\
RF-09 & El sistema deberá ofrecer una agenda digital por día y bahía,
donde las citas solicitadas en línea quedan ``Pendientes'' hasta que la
recepcionista verifique disponibilidad y las confirme manualmente & M2
-- Citas y Solicitudes & Should & HU-03, HU-04, HU-06 \\
RF-10 & El sistema deberá impedir la entrega física del vehículo con
saldo pendiente, a menos que el gerente autorice la entrega a crédito
para vehículos de flotillas empresariales previamente registradas. & M7
-- Portal del Cliente, Cobro y Gerencia & Must & HU-21, HU-22, HU-23 \\
RF-11 & El sistema deberá ofrecer al gerente un panel con las órdenes
abiertas y su estado, la ocupación de las 4 bahías, los tiempos de
estancia y los ingresos acumulados. & M7 -- Portal del Cliente, Cobro y
Gerencia & Should & HU-29 \\
\end{longtable}
}

\subsection{4.5 Historias de usuario (Product
Backlog)}\label{historias-de-usuario-product-backlog}

{\def\LTcaptype{none} % do not increment counter
\begin{longtable}[]{@{}
  >{\raggedright\arraybackslash}p{(\linewidth - 10\tabcolsep) * \real{0.1389}}
  >{\raggedright\arraybackslash}p{(\linewidth - 10\tabcolsep) * \real{0.4124}}
  >{\raggedright\arraybackslash}p{(\linewidth - 10\tabcolsep) * \real{0.1068}}
  >{\raggedright\arraybackslash}p{(\linewidth - 10\tabcolsep) * \real{0.0855}}
  >{\raggedright\arraybackslash}p{(\linewidth - 10\tabcolsep) * \real{0.1282}}
  >{\raggedright\arraybackslash}p{(\linewidth - 10\tabcolsep) * \real{0.1282}}@{}}
\toprule\noalign{}
\begin{minipage}[b]{\linewidth}\raggedright
\textbf{ID / issue}
\end{minipage} & \begin{minipage}[b]{\linewidth}\raggedright
\textbf{Historia (Como\ldots{} quiero\ldots{} para\ldots)}
\end{minipage} & \begin{minipage}[b]{\linewidth}\raggedright
\textbf{Prioridad}
\end{minipage} & \begin{minipage}[b]{\linewidth}\raggedright
\textbf{Puntos}
\end{minipage} & \begin{minipage}[b]{\linewidth}\raggedright
\textbf{Release}
\end{minipage} & \begin{minipage}[b]{\linewidth}\raggedright
\textbf{Estado}
\end{minipage} \\
\midrule\noalign{}
\endhead
\bottomrule\noalign{}
\endlastfoot
HU-01 / \#31 & Como recepcionista del taller, quiero registrar la
recepción del vehículo adjuntando kilometraje, nivel de combustible,
checklist de carrocería y fotografías del estado inicial, para contar
con un respaldo digital inalterable que evite reclamos infundados por
daños o faltantes al entregar la unidad. & Must & 3 & v0.1.0 &
Pendiente \\
HU-02 / \#32 & Como cliente del taller, quiero recibir un enlace
web seguro con la cotización detallada (mano de obra y refacciones) y
botones claros de aprobación o rechazo, para revisar y autorizar los
trabajos desde mi teléfono sin necesidad de realizar llamadas ni acudir
físicamente al taller & Must & 3 & v0.1.0 & Pendiente \\
HU-03 / \#33 & Como cliente del taller, quiero solicitar una cita
en línea seleccionando el día y la hora disponible, para apartar un
espacio sin necesidad de realizar llamadas telefónicas. & Must & 3 &
v0.1.0 & Pendiente \\
HU-04 / \#34 & Como recepcionista del taller, quiero consultar la
agenda digital por día y por bahía, para verificar las solicitudes
pendientes de los clientes y realizar su confirmación manual. & Must & 3
& v0.1.0 & Pendiente \\
HU-05 / \#35 & Como recepcionista del taller, quiero buscar a un
cliente o vehículo registrado mediante su número telefónico o placas,
para agilizar el proceso de recepción en visitas recurrentes. & Should &
2 & Fase 2 & Pendiente \\
HU-06 / \#36 & Como recepcionista del taller, quiero registrar un
vehículo con cita previa seleccionando su folio de origen, para que el
sistema precargue automáticamente los datos del cliente y del auto. &
Should & 2 & v0.1.0 & Pendiente \\
HU-07 / \#37 & Como recepcionista del taller, quiero ver la alerta
de "Sin Bahías Disponibles" cuando las 4 bahías físicas estén ocupadas,
para ubicar temporalmente el vehículo en el estacionamiento de
recepción. & Must & 3 & v0.1.0 & Pendiente \\
HU-08 / \#38 & Como jefe de mecánicos, quiero asignar la orden de
servicio a una de las 4 bahías físicas disponibles y a un mecánico
responsable, para iniciar formalmente los trabajos en el taller. & Must
& 5 & v0.1.0 & Pendiente \\
HU-09 / \#39 & Como jefe de mecánicos, quiero registrar el
diagnóstico técnico y las observaciones iniciales en la tableta, para
cambiar el estado de la orden a "En Diagnóstico". & Must & 3 & v0.1.0 &
Pendiente \\
HU-10 / \#40 & Como jefe de mecánicos, quiero registrar las
refacciones e insumos que los mecánicos detecten durante la revisión,
para que el área de almacén gestione su cotización o pedido. & Should &
3 & v0.1.0 & Pendiente \\
HU-11 / \#41 & Como jefe de mecánicos, quiero visualizar el
estatus actual de las 4 bahías de trabajo en una interfaz simplificada
para tabletas, para conocer rápidamente qué unidades están listas o en
proceso. & Should & 3 & Fase 2 & Pendiente \\
HU-12 / \#42 & Como recepcionista o jefe de mecánicos, quiero
generar un presupuesto desglosado de mano de obra y refacciones con
vigencia de 5 días naturales, para enviarlo de forma transparente al
cliente. & Must & 5 & v0.1.0 & Pendiente \\
HU-13 / \#43 & Como cliente del taller, quiero seleccionar o
deseleccionar servicios opcionales dentro de mi cotización antes de
aprobarla, para adaptar el presupuesto a mis necesidades financieras. &
Should & 3 & v0.1.0 & Pendiente \\
HU-14 / \#44 & Como cliente del taller, quiero realizar la firma
digital de aprobación mediante un código de verificación de 2 pasos
enviado a mi correo, para autorizar los trabajos de forma segura ya sea
de forma presencial o remota. & Must & 5 & v0.1.0 & Pendiente \\
HU-15 / \#45 & Como cliente del taller, quiero presionar un botón
de rechazo en la cotización digital, para cancelar la orden y coordinar
la recolección del vehículo sin autorizar reparaciones. & Should & 2 &
v0.1.0 & Pendiente \\
HU-16 / \#46 & Como cliente del taller, quiero que se me pida una
nueva autorización con código 2FA cuando los costos imprevistos superen
el 10 \% del presupuesto autorizado, para no pagar trabajos que no
aprobé. & Must & 5 & v0.1.0 & Pendiente \\
HU-17 / \#47 & Como encargado de almacén, quiero consultar y
administrar un directorio de proveedores locales con teléfonos y correos
electrónicos, para solicitar refacciones especiales rápidamente. &
Should & 3 & v0.1.0 & Pendiente \\
HU-18 / \#48 & Como encargado de almacén, quiero registrar los
costos, fechas de llegada y tickets de compra de las refacciones
asociadas a una orden, para mantener un control financiero estricto. &
Should & 3 & v0.1.0 & Pendiente \\
HU-19 / \#49 & Como jefe de mecánicos, quiero verificar el stock
de refacciones comunes en el inventario interno, para saber si se
dispone de la pieza de forma inmediata. & Could & 2 & Fase 2 &
Pendiente \\
HU-20 / \#50 & Como recepcionista, quiero recibir una notificación
automática cuando los trabajos del vehículo finalicen, para avisar al
cliente que su unidad está lista para entrega. & Should & 2 & v0.1.0 &
Pendiente \\
HU-21 / \#51 & Como gerente del taller, quiero que el sistema
impida registrar la entrega de un vehículo cuya orden tenga saldo
pendiente, para evitar pérdidas financieras al taller. & Must & 3 &
v0.1.0 & Pendiente \\
HU-22 / \#52 & Como gerente del taller, quiero autorizar de forma
excepcional la entrega a crédito únicamente para vehículos de flotillas
empresariales registradas, para cumplir con los convenios comerciales. &
Should & 3 & v0.1.0 & Pendiente \\
HU-23 / \#53 & Como recepcionista, quiero registrar el pago total
del servicio y cerrar formalmente la orden, para liberar la bahía de
trabajo y archivar el historial del vehículo. & Must & 3 & v0.1.0 &
Pendiente \\
HU-24 / \#54 & Como cliente, quiero consultar la vigencia de la
garantía de 30 días en mano de obra asociada a mi servicio anterior,
para reclamar ajustes en caso de incidencias. & Could & 2 & Fase 2 &
Pendiente \\
HU-25 / \#55 & Como recepcionista o jefe de mecánicos, quiero
registrar la llegada de autos, tomar fotografías y hacer diagnósticos
sin conexión a internet usando la batería del equipo, para que la
operación no se detenga ante fallas eléctricas. & Must & 8 & v0.1.0 &
Pendiente \\
HU-26 / \#56 & Como recepcionista o jefe de mecánicos, quiero que
los registros y fotografías que capturé sin conexión se sincronicen
automáticamente cuando vuelva la red, para no perder ni duplicar
información. & Must & 8 & v0.1.0 & Pendiente \\
HU-27 / \#57 & Como gerente del taller, quiero dar de alta al
personal y asignarle su rol (Gerente, Recepcionista, Jefe de Mecánicos o
Encargado de Almacén), para que cada quien acceda solo a la información
que le corresponde. & Must & 3 & v0.1.0 & Pendiente \\
HU-28 / \#58 & Como cliente, quiero visualizar la plataforma web
desde mi teléfono celular de forma responsiva y sin instalar
aplicaciones adicionales, para consultar mi estatus con facilidad. &
Must & 5 & v0.1.0 & Pendiente \\
HU-29 / \#59 & Como gerente, quiero generar reportes mensuales de
ingresos, servicios realizados y tiempos de ocupación de las 4 bahías,
para evaluar el rendimiento general del taller. & Could & 5 & Fase 2 &
Pendiente \\
HU-30 / \#60 & Como cliente del taller, quiero conocer el aviso de
privacidad antes de proporcionar mis datos, para saber para qué se
usarán conforme a la LFPDPPP. & Must & 3 & v0.1.0 & Pendiente \\
\end{longtable}
}

Detalle de las historias de usuario

\textbf{HU-01: Recepción e Inspección con Evidencia Fotográfica}

{\def\LTcaptype{none} % do not increment counter
\begin{longtable}[]{@{}
  >{\raggedright\arraybackslash}p{(\linewidth - 2\tabcolsep) * \real{0.2778}}
  >{\raggedright\arraybackslash}p{(\linewidth - 2\tabcolsep) * \real{0.7222}}@{}}
\toprule\noalign{}
\begin{minipage}[b]{\linewidth}\raggedright
\textbf{Campo}
\end{minipage} & \begin{minipage}[b]{\linewidth}\raggedright
\textbf{Contenido}
\end{minipage} \\
\midrule\noalign{}
\endhead
\bottomrule\noalign{}
\endlastfoot
ID e issue & HU-01 \\
Historia & Como recepcionista del taller, quiero registrar la recepción
del vehículo adjuntando kilometraje, nivel de combustible, checklist de
carrocería y fotografías del estado inicial, para contar con un respaldo
digital inalterable que evite reclamos infundados por daños o faltantes
al entregar la unidad. \\
Criterio de aceptación 1 & Dado que la recepcionista se encuentra en el
módulo de recepción vehicular, cuando intenta registrar una unidad
nueva, entonces el sistema exige obligatoriamente adjuntar al menos 4
fotografías (frente, parte trasera, costado izquierdo y costado
derecho). \\
Criterio de aceptación 2 & Dado que se capturan las evidencias, cuando
se finaliza el formulario de inspección, entonces el sistema genera una
firma visual o sello digital con fecha y hora exacta de recepción,
bloqueando el avance al estado "En Diagnóstico" si falta completar algún
punto del checklist obligatorio. \\
Reglas de negocio & RN-12 (Evidencia fotográfica de recepción), RN-07
(Clasificación de vehículos autorizados). \\
Wireframe & W-01 (4.10, Review 2) \\
Casos de prueba & CPA-01 a CPA-04 \\
\end{longtable}
}

\textbf{HU-02: Autorización Digital de Cotización por Parte del Cliente}

{\def\LTcaptype{none} % do not increment counter
\begin{longtable}[]{@{}
  >{\raggedright\arraybackslash}p{(\linewidth - 2\tabcolsep) * \real{0.2778}}
  >{\raggedright\arraybackslash}p{(\linewidth - 2\tabcolsep) * \real{0.7222}}@{}}
\toprule\noalign{}
\begin{minipage}[b]{\linewidth}\raggedright
\textbf{Campo}
\end{minipage} & \begin{minipage}[b]{\linewidth}\raggedright
\textbf{Contenido}
\end{minipage} \\
\midrule\noalign{}
\endhead
\bottomrule\noalign{}
\endlastfoot
ID e issue & HU-02 \\
Historia & Como cliente del taller, quiero recibir un enlace web seguro
con la cotización detallada (mano de obra y refacciones) y botones
claros de aprobación o rechazo, para revisar y autorizar los trabajos
desde mi teléfono sin necesidad de realizar llamadas ni acudir
físicamente al taller \\
Criterio de aceptación 1 & Dado que la cotización ha sido generada por
el taller, cuando el sistema envía el aviso, entonces el cliente recibe
un enlace por correo electrónico (o por SMS cuando se habilite RF-02.3)
que no requiere registro previo ni contraseña, operando mediante un
token temporal firmado. \\
Criterio de aceptación 2 & Dado que el cliente revisa el desglose en su
dispositivo móvil, cuando presiona el botón "Autorizar" e ingresa un
código 2FA válido (HU-14), entonces el sistema registra de forma
inalterable la dirección IP, fecha, hora y el monto final aceptado,
cambiando automáticamente el estado de la orden a "Autorizada" y
notificando al mecánico responsable. \\
Reglas de negocio & RN-01 (Expiración y vigencia de enlaces), RN-02
(Doble verificación en autorizaciones), RN-11 (Vigencia de
cotizaciones). \\
Wireframe & W-04 (4.10, Review 2) \\
Casos de prueba & CPA-10, CPA-30 \\
\end{longtable}
}

\textbf{HU-03: Solicitud de Cita en Línea}

{\def\LTcaptype{none} % do not increment counter
\begin{longtable}[]{@{}
  >{\raggedright\arraybackslash}p{(\linewidth - 2\tabcolsep) * \real{0.2778}}
  >{\raggedright\arraybackslash}p{(\linewidth - 2\tabcolsep) * \real{0.7222}}@{}}
\toprule\noalign{}
\begin{minipage}[b]{\linewidth}\raggedright
\textbf{Campo}
\end{minipage} & \begin{minipage}[b]{\linewidth}\raggedright
\textbf{Contenido}
\end{minipage} \\
\midrule\noalign{}
\endhead
\bottomrule\noalign{}
\endlastfoot
ID e issue & HU-03 \\
Historia & Como cliente del taller, quiero solicitar una cita en línea
seleccionando el día y la hora disponible, para apartar un espacio sin
necesidad de realizar llamadas telefónicas. \\
Criterio de aceptación 1 & Dado que el cliente accede al portal web,
cuando selecciona fecha, hora y tipo de servicio, entonces el sistema
valida que exista disponibilidad en la agenda del taller. \\
Criterio de aceptación 2 & Dado que se envía la solicitud, cuando se
registra con éxito, entonces queda en estado "Pendiente" a la espera de
la confirmación manual de la recepcionista. \\
Reglas de negocio & RN-08 (Confirmación manual de citas web). \\
Wireframe & W-05 (4.10, Review 2) \\
Casos de prueba & CPA-31, CPA-34 \\
\end{longtable}
}

\textbf{HU-04: Consulta y Confirmación Manual de Citas}

{\def\LTcaptype{none} % do not increment counter
\begin{longtable}[]{@{}
  >{\raggedright\arraybackslash}p{(\linewidth - 2\tabcolsep) * \real{0.2778}}
  >{\raggedright\arraybackslash}p{(\linewidth - 2\tabcolsep) * \real{0.7222}}@{}}
\toprule\noalign{}
\begin{minipage}[b]{\linewidth}\raggedright
\textbf{Campo}
\end{minipage} & \begin{minipage}[b]{\linewidth}\raggedright
\textbf{Contenido}
\end{minipage} \\
\midrule\noalign{}
\endhead
\bottomrule\noalign{}
\endlastfoot
ID e issue & HU-04 \\
Historia & Como recepcionista del taller, quiero consultar la agenda
digital por día y por bahía, para verificar las solicitudes pendientes
de los clientes y realizar su confirmación manual. \\
Criterio de aceptación 1 & Dado que la recepcionista accede al módulo de
agenda, cuando filtra los registros por el día actual, entonces el
sistema despliega todas las citas con estatus "Pendiente". \\
Criterio de aceptación 2 & Dado que la recepcionista revisa una cita
solicitada, cuando presiona el botón "Confirmar", entonces el estatus
cambia a "Confirmada" y se habilita la asignación de la bahía
correspondiente. \\
Reglas de negocio & RN-08 (Confirmación manual de citas web), RN-05
(Capacidad física del taller). \\
Wireframe & W-05 (4.10, Review 2) \\
Casos de prueba & CPA-32 \\
\end{longtable}
}

\textbf{HU-05: Búsqueda Rápida de Clientes y Vehículos}

{\def\LTcaptype{none} % do not increment counter
\begin{longtable}[]{@{}
  >{\raggedright\arraybackslash}p{(\linewidth - 2\tabcolsep) * \real{0.2778}}
  >{\raggedright\arraybackslash}p{(\linewidth - 2\tabcolsep) * \real{0.7222}}@{}}
\toprule\noalign{}
\begin{minipage}[b]{\linewidth}\raggedright
\textbf{Campo}
\end{minipage} & \begin{minipage}[b]{\linewidth}\raggedright
\textbf{Contenido}
\end{minipage} \\
\midrule\noalign{}
\endhead
\bottomrule\noalign{}
\endlastfoot
ID e issue & HU-05 \\
Historia & Como recepcionista del taller, quiero buscar a un cliente o
vehículo registrado mediante su número telefónico o placas, para
agilizar el proceso de recepción en visitas recurrentes. \\
Criterio de aceptación 1 & Dado que la recepcionista se encuentra en el
módulo de búsqueda, cuando ingresa los dígitos de las placas o el
teléfono, entonces el sistema filtra instantáneamente la base de datos
de clientes recurrentes. \\
Criterio de aceptación 2 & Dado que se localiza el expediente, cuando el
usuario lo selecciona, entonces el sistema precarga los datos generales
del cliente y del auto para evitar duplicidad de registros. \\
Reglas de negocio & Sin regla de negocio específica; ver RF-03.1. \\
Wireframe & Se define en 4.10 (Review 2) \\
Casos de prueba & Se definen en 5.2 (Review 2) \\
\end{longtable}
}

\textbf{HU-06: Registro de Vehículos con Cita Previa}

{\def\LTcaptype{none} % do not increment counter
\begin{longtable}[]{@{}
  >{\raggedright\arraybackslash}p{(\linewidth - 2\tabcolsep) * \real{0.2778}}
  >{\raggedright\arraybackslash}p{(\linewidth - 2\tabcolsep) * \real{0.7222}}@{}}
\toprule\noalign{}
\begin{minipage}[b]{\linewidth}\raggedright
\textbf{Campo}
\end{minipage} & \begin{minipage}[b]{\linewidth}\raggedright
\textbf{Contenido}
\end{minipage} \\
\midrule\noalign{}
\endhead
\bottomrule\noalign{}
\endlastfoot
ID e issue & HU-06 \\
Historia & Como recepcionista del taller, quiero registrar un vehículo
con cita previa seleccionando su folio de origen, para que el sistema
precargue automáticamente los datos del cliente y del auto. \\
Criterio de aceptación 1 & Dado que el cliente llega al taller con cita
agendada, cuando la recepcionista introduce el número de folio de la
cita, entonces el sistema recupera la información asociada. \\
Criterio de aceptación 2 & Dado que los datos están precargados, cuando
se procede con el ingreso, entonces el sistema vincula la cita
directamente con la nueva orden de servicio abierta. \\
Reglas de negocio & RN-08 (Confirmación manual de citas web). \\
Wireframe & W-01 (4.10, Review 2) \\
Casos de prueba & CPA-33 \\
\end{longtable}
}

\textbf{HU-07: Alerta de Bahías Ocupadas}

{\def\LTcaptype{none} % do not increment counter
\begin{longtable}[]{@{}
  >{\raggedright\arraybackslash}p{(\linewidth - 2\tabcolsep) * \real{0.2778}}
  >{\raggedright\arraybackslash}p{(\linewidth - 2\tabcolsep) * \real{0.7222}}@{}}
\toprule\noalign{}
\begin{minipage}[b]{\linewidth}\raggedright
\textbf{Campo}
\end{minipage} & \begin{minipage}[b]{\linewidth}\raggedright
\textbf{Contenido}
\end{minipage} \\
\midrule\noalign{}
\endhead
\bottomrule\noalign{}
\endlastfoot
ID e issue & HU-07 \\
Historia & Como recepcionista del taller, quiero ver la alerta de "Sin
Bahías Disponibles" cuando las 4 bahías físicas estén ocupadas, para
ubicar temporalmente el vehículo en el estacionamiento de recepción. \\
Criterio de aceptación 1 & Dado que las 4 bahías físicas del taller
registran una orden activa, cuando la recepcionista intenta registrar la
entrada de un nuevo vehículo, entonces el sistema emite una advertencia
visual en pantalla indicando "Sin Bahías Disponibles". \\
Criterio de aceptación 2 & Dado que se activa la alerta de saturación,
cuando el usuario acepta la notificación, entonces el sistema permite
guardar el registro temporal asignando la unidad al estacionamiento de
espera. \\
Reglas de negocio & RN-05 (Capacidad física del taller). \\
Wireframe & W-01 (4.10, Review 2) \\
Casos de prueba & CPA-05, CPA-06 \\
\end{longtable}
}

\textbf{HU-08: Asignación de Bahía y Mecánico}

{\def\LTcaptype{none} % do not increment counter
\begin{longtable}[]{@{}
  >{\raggedright\arraybackslash}p{(\linewidth - 2\tabcolsep) * \real{0.2778}}
  >{\raggedright\arraybackslash}p{(\linewidth - 2\tabcolsep) * \real{0.7222}}@{}}
\toprule\noalign{}
\begin{minipage}[b]{\linewidth}\raggedright
\textbf{Campo}
\end{minipage} & \begin{minipage}[b]{\linewidth}\raggedright
\textbf{Contenido}
\end{minipage} \\
\midrule\noalign{}
\endhead
\bottomrule\noalign{}
\endlastfoot
ID e issue & HU-08 \\
Historia & Como jefe de mecánicos, quiero asignar la orden de servicio a
una de las 4 bahías físicas disponibles y a un mecánico responsable,
para iniciar formalmente los trabajos en el taller. \\
Criterio de aceptación 1 & Dado que la orden se encuentra abierta,
cuando el jefe de mecánicos selecciona una de las 4 bahías físicas
libres, entonces el sistema actualiza el estado de la bahía a
"Ocupada". \\
Criterio de aceptación 2 & Dado que las 4 bahías están ocupadas, cuando
se intenta una asignación, entonces el sistema emite una alerta
impidiendo el cruce de unidades en la misma bahía. \\
Reglas de negocio & RN-05 (Capacidad física del taller), RN-09
(Asignación y colaboración de mecánicos). \\
Wireframe & W-02 (4.10, Review 2) \\
Casos de prueba & CPA-07, CPA-08 \\
\end{longtable}
}

\textbf{HU-09: Registro de Diagnóstico en Tableta}

{\def\LTcaptype{none} % do not increment counter
\begin{longtable}[]{@{}
  >{\raggedright\arraybackslash}p{(\linewidth - 2\tabcolsep) * \real{0.2778}}
  >{\raggedright\arraybackslash}p{(\linewidth - 2\tabcolsep) * \real{0.7222}}@{}}
\toprule\noalign{}
\begin{minipage}[b]{\linewidth}\raggedright
\textbf{Campo}
\end{minipage} & \begin{minipage}[b]{\linewidth}\raggedright
\textbf{Contenido}
\end{minipage} \\
\midrule\noalign{}
\endhead
\bottomrule\noalign{}
\endlastfoot
ID e issue & HU-09 \\
Historia & Como jefe de mecánicos, quiero registrar el diagnóstico
técnico y las observaciones iniciales en la tableta, para cambiar el
estado de la orden a "En Diagnóstico". \\
Criterio de aceptación 1 & Dado que el jefe de mecánicos se encuentra
autenticado en la tableta, cuando selecciona una orden con vehículo en
bahía, entonces visualiza un formulario táctil optimizado con campos
amplios para el dictamen técnico. \\
Criterio de aceptación 2 & Dado que se ingresan las observaciones de la
revisión, cuando se presiona "Guardar Diagnóstico", entonces el sistema
actualiza automáticamente el estado de la orden a "En Diagnóstico" y
habilita la creación de cotizaciones. \\
Reglas de negocio & RN-12 (Evidencia fotográfica de recepción).
Usabilidad en tableta: RNF-02 y RNF-03. \\
Wireframe & W-02 (4.10, Review 2) \\
Casos de prueba & CPA-02, CPA-03 \\
\end{longtable}
}

\textbf{HU-10: Solicitud de Refacciones e Insumos}

{\def\LTcaptype{none} % do not increment counter
\begin{longtable}[]{@{}
  >{\raggedright\arraybackslash}p{(\linewidth - 2\tabcolsep) * \real{0.2778}}
  >{\raggedright\arraybackslash}p{(\linewidth - 2\tabcolsep) * \real{0.7222}}@{}}
\toprule\noalign{}
\begin{minipage}[b]{\linewidth}\raggedright
\textbf{Campo}
\end{minipage} & \begin{minipage}[b]{\linewidth}\raggedright
\textbf{Contenido}
\end{minipage} \\
\midrule\noalign{}
\endhead
\bottomrule\noalign{}
\endlastfoot
ID e issue & HU-10 \\
Historia & Como jefe de mecánicos, quiero registrar las refacciones e
insumos que los mecánicos detecten durante la revisión, para que el área
de almacén gestione su cotización o pedido. \\
Criterio de aceptación 1 & Dado que un mecánico reporta piezas
necesarias en la bahía, cuando el jefe de mecánicos las añade desde la
tableta y las vincula a la orden, entonces la solicitud se envía
automáticamente al módulo de almacén. \\
Criterio de aceptación 2 & Dado que el almacén recibe el requerimiento,
cuando se procesa, entonces el sistema refleja el estatus de la
refacción como "Solicitada" o "Pedida". \\
Reglas de negocio & RN-04 (Re-cotización obligatoria por sobrecostos),
RN-06 (Reubicación por espera de refacciones). \\
Wireframe & Se define en 4.10 (Review 2) \\
Casos de prueba & Se definen en 5.2 (Review 2) \\
\end{longtable}
}

\textbf{HU-11: Visualización del Estatus de Bahías}

{\def\LTcaptype{none} % do not increment counter
\begin{longtable}[]{@{}
  >{\raggedright\arraybackslash}p{(\linewidth - 2\tabcolsep) * \real{0.2778}}
  >{\raggedright\arraybackslash}p{(\linewidth - 2\tabcolsep) * \real{0.7222}}@{}}
\toprule\noalign{}
\begin{minipage}[b]{\linewidth}\raggedright
\textbf{Campo}
\end{minipage} & \begin{minipage}[b]{\linewidth}\raggedright
\textbf{Contenido}
\end{minipage} \\
\midrule\noalign{}
\endhead
\bottomrule\noalign{}
\endlastfoot
ID e issue & HU-11 \\
Historia & Como jefe de mecánicos, quiero visualizar el estatus actual
de las 4 bahías de trabajo en una interfaz simplificada para tabletas,
para conocer rápidamente qué unidades están listas o en proceso. \\
Criterio de aceptación 1 & Dado que el jefe de mecánicos accede al
tablero de control, cuando visualiza la sección de bahías, entonces
identifica mediante códigos de color las unidades activas, en
diagnóstico o listas. \\
Criterio de aceptación 2 & Dado que se actualiza el estado de una
reparación, cuando cambia de fase, entonces el tablero refresca la
información en tiempo real para optimizar la supervisión del taller. \\
Reglas de negocio & RN-05 (Capacidad física del taller). Usabilidad en
tableta: RNF-02. \\
Wireframe & W-02 (4.10, Review 2) \\
Casos de prueba & Se definen en 5.2 (Review 2) \\
\end{longtable}
}

\textbf{HU-12: Generación de Presupuesto Desglosado}

{\def\LTcaptype{none} % do not increment counter
\begin{longtable}[]{@{}
  >{\raggedright\arraybackslash}p{(\linewidth - 2\tabcolsep) * \real{0.2778}}
  >{\raggedright\arraybackslash}p{(\linewidth - 2\tabcolsep) * \real{0.7222}}@{}}
\toprule\noalign{}
\begin{minipage}[b]{\linewidth}\raggedright
\textbf{Campo}
\end{minipage} & \begin{minipage}[b]{\linewidth}\raggedright
\textbf{Contenido}
\end{minipage} \\
\midrule\noalign{}
\endhead
\bottomrule\noalign{}
\endlastfoot
ID e issue & HU-12 \\
Historia & Como recepcionista o jefe de mecánicos, quiero generar un
presupuesto desglosado de mano de obra y refacciones con vigencia de 5
días naturales, para enviarlo de forma transparente al cliente. \\
Criterio de aceptación 1 & Dado que el diagnóstico ha concluido, cuando
el personal selecciona la opción de cotizar, entonces el sistema
presenta apartados separados para costos de mano de obra y
refacciones. \\
Criterio de aceptación 2 & Dado que se emite la cotización, entonces el
sistema calcula automáticamente el total con IVA e integra una fecha de
expiración exactamente a los 5 días naturales posteriores a su
emisión. \\
Reglas de negocio & RN-11 (Vigencia de cotizaciones), RN-03 (Selección
granular de rubros). \\
Wireframe & W-03 (4.10, Review 2) \\
Casos de prueba & CPA-10 \\
\end{longtable}
}

\textbf{HU-13: Selección de Servicios Opcionales en Cotización}

{\def\LTcaptype{none} % do not increment counter
\begin{longtable}[]{@{}
  >{\raggedright\arraybackslash}p{(\linewidth - 2\tabcolsep) * \real{0.2778}}
  >{\raggedright\arraybackslash}p{(\linewidth - 2\tabcolsep) * \real{0.7222}}@{}}
\toprule\noalign{}
\begin{minipage}[b]{\linewidth}\raggedright
\textbf{Campo}
\end{minipage} & \begin{minipage}[b]{\linewidth}\raggedright
\textbf{Contenido}
\end{minipage} \\
\midrule\noalign{}
\endhead
\bottomrule\noalign{}
\endlastfoot
ID e issue & HU-13 \\
Historia & Como cliente del taller, quiero seleccionar o deseleccionar
servicios opcionales dentro de mi cotización antes de aprobarla, para
adaptar el presupuesto a mis necesidades financieras. \\
Criterio de aceptación 1 & Dado que el cliente revisa el presupuesto
desglosado en su celular, cuando marca o desmarca las casillas de
servicios opcionales o sugeridos, entonces el sistema recalcula el costo
total en tiempo real. \\
Criterio de aceptación 2 & Dado que se ajustan los opcionales, cuando el
cliente procede a la firma, entonces el sistema guarda únicamente los
conceptos seleccionados como parte oficial de la orden. \\
Reglas de negocio & RN-03 (Selección granular de rubros). \\
Wireframe & W-04 (4.10, Review 2) \\
Casos de prueba & CPA-12 \\
\end{longtable}
}

\textbf{HU-14: Firma Digital y Autorización por 2FA}

{\def\LTcaptype{none} % do not increment counter
\begin{longtable}[]{@{}
  >{\raggedright\arraybackslash}p{(\linewidth - 2\tabcolsep) * \real{0.2778}}
  >{\raggedright\arraybackslash}p{(\linewidth - 2\tabcolsep) * \real{0.7222}}@{}}
\toprule\noalign{}
\begin{minipage}[b]{\linewidth}\raggedright
\textbf{Campo}
\end{minipage} & \begin{minipage}[b]{\linewidth}\raggedright
\textbf{Contenido}
\end{minipage} \\
\midrule\noalign{}
\endhead
\bottomrule\noalign{}
\endlastfoot
ID e issue & HU-14 \\
Historia & Como cliente del taller, quiero realizar la firma digital de
aprobación mediante un código de verificación de 2 pasos enviado a mi
correo, para autorizar los trabajos de forma segura ya sea de forma
presencial o remota. \\
Criterio de aceptación 1 & Dado que el cliente accede al enlace web de
cotización, when solicita el código de autorización, entonces el sistema
envía un token numérico de 6 dígitos al correo electrónico registrado
del cliente. \\
Criterio de aceptación 2 & Dado que el cliente ingresa el código
recibido en la interfaz, cuando presiona "Autorizar Presupuesto",
entonces el sistema valida el token, registra la dirección IP y la marca
de tiempo, cambiando la orden a "Autorizada". \\
Reglas de negocio & RN-02 (Doble verificación en autorizaciones). No
repudio: RNF-04. \\
Wireframe & W-04 (4.10, Review 2) \\
Casos de prueba & CPA-11, CPA-13 a CPA-16 \\
\end{longtable}
}

\textbf{HU-15: Rechazo de Cotización y Cancelación}

{\def\LTcaptype{none} % do not increment counter
\begin{longtable}[]{@{}
  >{\raggedright\arraybackslash}p{(\linewidth - 2\tabcolsep) * \real{0.2778}}
  >{\raggedright\arraybackslash}p{(\linewidth - 2\tabcolsep) * \real{0.7222}}@{}}
\toprule\noalign{}
\begin{minipage}[b]{\linewidth}\raggedright
\textbf{Campo}
\end{minipage} & \begin{minipage}[b]{\linewidth}\raggedright
\textbf{Contenido}
\end{minipage} \\
\midrule\noalign{}
\endhead
\bottomrule\noalign{}
\endlastfoot
ID e issue & HU-15 \\
Historia & Como cliente del taller, quiero presionar un botón de rechazo
en la cotización digital, para cancelar la orden y coordinar la
recolección del vehículo sin autorizar reparaciones. \\
Criterio de aceptación 1 & Dado que el cliente evalúa la cotización en
su dispositivo, cuando presiona el botón "Rechazar", entonces el sistema
solicita una confirmación de la acción. \\
Criterio de aceptación 2 & Dado que se confirma el rechazo, cuando se
procesa, entonces la orden cambia a estatus "Rechazada" y notifica a
recepción para coordinar la entrega física del auto sin costo de
reparación. \\
Reglas de negocio & RN-02 (Doble verificación en autorizaciones). Flujo
de rechazo: CU-11 (A1). \\
Wireframe & W-04 (4.10, Review 2) \\
Casos de prueba & CPA-17 \\
\end{longtable}
}

\textbf{HU-16: Re-cotización por Incremento \textgreater10\%}

{\def\LTcaptype{none} % do not increment counter
\begin{longtable}[]{@{}
  >{\raggedright\arraybackslash}p{(\linewidth - 2\tabcolsep) * \real{0.2778}}
  >{\raggedright\arraybackslash}p{(\linewidth - 2\tabcolsep) * \real{0.7222}}@{}}
\toprule\noalign{}
\begin{minipage}[b]{\linewidth}\raggedright
\textbf{Campo}
\end{minipage} & \begin{minipage}[b]{\linewidth}\raggedright
\textbf{Contenido}
\end{minipage} \\
\midrule\noalign{}
\endhead
\bottomrule\noalign{}
\endlastfoot
ID e issue & HU-16 \\
Historia & Como cliente del taller, quiero que se me pida una nueva
autorización con código 2FA cuando los costos imprevistos superen el 10
\% del presupuesto autorizado, para no pagar trabajos que no aprobé. \\
Criterio de aceptación 1 & Dado que surgen refacciones imprevistas,
cuando el incremento total supera el 10\% del presupuesto aprobado,
entonces el sistema bloquea el flujo y solicita una nueva autorización
con código 2FA. \\
Criterio de aceptación 2 & Dado que el cliente recibe la notificación
del sobrecosto, cuando aprueba mediante el código digital, entonces se
actualiza el monto total e historial de la orden. \\
Reglas de negocio & RN-04 (Re-cotización obligatoria por sobrecostos),
RN-02 (Doble verificación en autorizaciones). \\
Wireframe & W-03 (4.10, Review 2) \\
Casos de prueba & CPA-18, CPA-19 \\
\end{longtable}
}

\textbf{HU-17: Directorio y Gestión de Proveedores Locales}

{\def\LTcaptype{none} % do not increment counter
\begin{longtable}[]{@{}
  >{\raggedright\arraybackslash}p{(\linewidth - 2\tabcolsep) * \real{0.2778}}
  >{\raggedright\arraybackslash}p{(\linewidth - 2\tabcolsep) * \real{0.7222}}@{}}
\toprule\noalign{}
\begin{minipage}[b]{\linewidth}\raggedright
\textbf{Campo}
\end{minipage} & \begin{minipage}[b]{\linewidth}\raggedright
\textbf{Contenido}
\end{minipage} \\
\midrule\noalign{}
\endhead
\bottomrule\noalign{}
\endlastfoot
ID e issue & HU-17 \\
Historia & Como encargado de almacén, quiero consultar y administrar un
directorio de proveedores locales con teléfonos y correos electrónicos,
para solicitar refacciones especiales rápidamente. \\
Criterio de aceptación 1 & Dado que el almacenista ingresa al módulo de
refacciones, cuando abre el directorio de proveedores, entonces
visualiza la lista con datos de contacto actualizados (teléfonos y
correos). \\
Criterio de aceptación 2 & Dado que se requiere dar de alta un nuevo
proveedor, cuando el almacenista completa el formulario de registro,
entonces el sistema lo almacena para futuras solicitudes de piezas. \\
Reglas de negocio & RN-06 (Reubicación por espera de refacciones). \\
Wireframe & Se define en 4.10 (Review 2) \\
Casos de prueba & Se definen en 5.2 (Review 2) \\
\end{longtable}
}

\textbf{HU-18: Registro de Costos e Historial de Compras}

{\def\LTcaptype{none} % do not increment counter
\begin{longtable}[]{@{}
  >{\raggedright\arraybackslash}p{(\linewidth - 2\tabcolsep) * \real{0.2778}}
  >{\raggedright\arraybackslash}p{(\linewidth - 2\tabcolsep) * \real{0.7222}}@{}}
\toprule\noalign{}
\begin{minipage}[b]{\linewidth}\raggedright
\textbf{Campo}
\end{minipage} & \begin{minipage}[b]{\linewidth}\raggedright
\textbf{Contenido}
\end{minipage} \\
\midrule\noalign{}
\endhead
\bottomrule\noalign{}
\endlastfoot
ID e issue & HU-18 \\
Historia & Como encargado de almacén, quiero registrar los costos,
fechas de llegada y tickets de compra de las refacciones asociadas a una
orden, para mantener un control financiero estricto. \\
Criterio de aceptación 1 & Dado que se adquiere una pieza externa con un
proveedor, cuando el almacenista registra el ticket, costo unitario y
fecha estimada de entrega, entonces el sistema lo vincula a la orden de
servicio. \\
Criterio de aceptación 2 & Dado que se ingresa la compra, cuando se
consolida el registro, entonces el sistema actualiza el historial de
adquisiciones de la refacción para auditoría interna. \\
Reglas de negocio & Sin regla de negocio específica; ver RF-08.2 y
RF-08.3. \\
Wireframe & Se define en 4.10 (Review 2) \\
Casos de prueba & Se definen en 5.2 (Review 2) \\
\end{longtable}
}

\textbf{HU-19: Verificación de Stock Interno}

{\def\LTcaptype{none} % do not increment counter
\begin{longtable}[]{@{}
  >{\raggedright\arraybackslash}p{(\linewidth - 2\tabcolsep) * \real{0.2778}}
  >{\raggedright\arraybackslash}p{(\linewidth - 2\tabcolsep) * \real{0.7222}}@{}}
\toprule\noalign{}
\begin{minipage}[b]{\linewidth}\raggedright
\textbf{Campo}
\end{minipage} & \begin{minipage}[b]{\linewidth}\raggedright
\textbf{Contenido}
\end{minipage} \\
\midrule\noalign{}
\endhead
\bottomrule\noalign{}
\endlastfoot
ID e issue & HU-19 \\
Historia & Como jefe de mecánicos, quiero verificar el stock de
refacciones comunes en el inventario interno, para saber si se dispone
de la pieza de forma inmediata. \\
Criterio de aceptación 1 & Dado que el jefe de mecánicos busca un insumo
estándar, cuando consulta el inventario en el sistema, entonces
visualiza las existencias disponibles en almacén. \\
Criterio de aceptación 2 & Dado que hay existencia en stock, cuando se
selecciona la pieza, entonces el sistema descuenta automáticamente la
cantidad del inventario interno. \\
Reglas de negocio & Sin regla de negocio específica; ver RF-08.1. \\
Wireframe & Se define en 4.10 (Review 2) \\
Casos de prueba & Se definen en 5.2 (Review 2) \\
\end{longtable}
}

\textbf{HU-20: Notificación Automática de Fin de Trabajos}

{\def\LTcaptype{none} % do not increment counter
\begin{longtable}[]{@{}
  >{\raggedright\arraybackslash}p{(\linewidth - 2\tabcolsep) * \real{0.2778}}
  >{\raggedright\arraybackslash}p{(\linewidth - 2\tabcolsep) * \real{0.7222}}@{}}
\toprule\noalign{}
\begin{minipage}[b]{\linewidth}\raggedright
\textbf{Campo}
\end{minipage} & \begin{minipage}[b]{\linewidth}\raggedright
\textbf{Contenido}
\end{minipage} \\
\midrule\noalign{}
\endhead
\bottomrule\noalign{}
\endlastfoot
ID e issue & HU-20 \\
Historia & Como recepcionista, quiero recibir una notificación
automática cuando los trabajos del vehículo finalicen, para avisar al
cliente que su unidad está lista para entrega. \\
Criterio de aceptación 1 & Dado que el mecánico marca la orden como
terminada en bahía, cuando se actualiza el estatus, entonces el sistema
envía una alerta interna a la recepcionista. \\
Criterio de aceptación 2 & Dado que la recepción recibe la alerta,
cuando se verifica el estado, entonces se habilita el aviso automatizado
al cliente de que su vehículo está listo. \\
Reglas de negocio & Sin regla de negocio específica; ver CU-07. \\
Wireframe & W-06 (4.10, Review 2) \\
Casos de prueba & Se definen en 5.2 (Review 2) \\
\end{longtable}
}

\textbf{HU-21: Bloqueo de Entrega por Saldo Pendiente}

{\def\LTcaptype{none} % do not increment counter
\begin{longtable}[]{@{}
  >{\raggedright\arraybackslash}p{(\linewidth - 2\tabcolsep) * \real{0.2778}}
  >{\raggedright\arraybackslash}p{(\linewidth - 2\tabcolsep) * \real{0.7222}}@{}}
\toprule\noalign{}
\begin{minipage}[b]{\linewidth}\raggedright
\textbf{Campo}
\end{minipage} & \begin{minipage}[b]{\linewidth}\raggedright
\textbf{Contenido}
\end{minipage} \\
\midrule\noalign{}
\endhead
\bottomrule\noalign{}
\endlastfoot
ID e issue & HU-21 \\
Historia & Como gerente del taller, quiero que el sistema impida
registrar la entrega de un vehículo cuya orden tenga saldo pendiente,
para evitar pérdidas financieras al taller. \\
Criterio de aceptación 1 & Dado que un vehículo se encuentra en estado
"Lista para Entrega", cuando la recepcionista intenta registrar la
salida y el balance de pago es mayor a cero, entonces el sistema bloquea
la acción de liberación. \\
Criterio de aceptación 2 & Dado que se activa el bloqueo por saldo
pendiente, cuando el sistema detecta que cuenta con un convenio de
crédito autorizado, entonces permite omitir la restricción
exclusivamente bajo permisos del rol Gerente. \\
Reglas de negocio & RN-10 (Restricción de entrega y excepción a
flotillas). \\
Wireframe & W-06 (4.10, Review 2) \\
Casos de prueba & CPA-21 \\
\end{longtable}
}

\textbf{HU-22: Autorización Excepcional de Crédito a Flotillas}

{\def\LTcaptype{none} % do not increment counter
\begin{longtable}[]{@{}
  >{\raggedright\arraybackslash}p{(\linewidth - 2\tabcolsep) * \real{0.2778}}
  >{\raggedright\arraybackslash}p{(\linewidth - 2\tabcolsep) * \real{0.7222}}@{}}
\toprule\noalign{}
\begin{minipage}[b]{\linewidth}\raggedright
\textbf{Campo}
\end{minipage} & \begin{minipage}[b]{\linewidth}\raggedright
\textbf{Contenido}
\end{minipage} \\
\midrule\noalign{}
\endhead
\bottomrule\noalign{}
\endlastfoot
ID e issue & HU-22 \\
Historia & Como gerente del taller, quiero autorizar de forma
excepcional la entrega a crédito únicamente para vehículos de flotillas
empresariales registradas, para cumplir con los convenios
comerciales. \\
Criterio de aceptación 1 & Dado que un vehículo pertenece a una empresa
con cuenta de flotilla registrada, cuando se intenta liberar sin pago en
caja, entonces el sistema solicita la clave de autorización
gerencial. \\
Criterio de aceptación 2 & Dado que el gerente ingresa su contraseña de
aprobación, cuando valida el convenio comercial, entonces el sistema
permite liberar la unidad y factura el adeudo a crédito empresarial. \\
Reglas de negocio & RN-10 (Restricción de entrega y excepción a
flotillas). \\
Wireframe & W-06 (4.10, Review 2) \\
Casos de prueba & CPA-24, CPA-25 \\
\end{longtable}
}

\textbf{HU-23: Registro de Pago Total y Cierre de Orden}

{\def\LTcaptype{none} % do not increment counter
\begin{longtable}[]{@{}
  >{\raggedright\arraybackslash}p{(\linewidth - 2\tabcolsep) * \real{0.2778}}
  >{\raggedright\arraybackslash}p{(\linewidth - 2\tabcolsep) * \real{0.7222}}@{}}
\toprule\noalign{}
\begin{minipage}[b]{\linewidth}\raggedright
\textbf{Campo}
\end{minipage} & \begin{minipage}[b]{\linewidth}\raggedright
\textbf{Contenido}
\end{minipage} \\
\midrule\noalign{}
\endhead
\bottomrule\noalign{}
\endlastfoot
ID e issue & HU-23 \\
Historia & Como recepcionista, quiero registrar el pago total del
servicio y cerrar formalmente la orden, para liberar la bahía de trabajo
y archivar el historial del vehículo. \\
Criterio de aceptación 1 & Dado que el cliente cubre el monto total en
caja, cuando la recepcionista registra el método de pago y confirma la
transacción, entonces el sistema salda la cuenta a cero. \\
Criterio de aceptación 2 & Dado que el pago ha sido registrado con
éxito, cuando se ejecuta el cierre, entonces la orden cambia a estado
"Cerrada", liberando de inmediato la bahía física ocupada y generando el
comprobante digital. \\
Reglas de negocio & RN-10 (Restricción de entrega y excepción a
flotillas), RN-13 (Garantía de mano de obra), RN-01 (Expiración y
vigencia de enlaces). \\
Wireframe & W-06 (4.10, Review 2) \\
Casos de prueba & CPA-22, CPA-23 \\
\end{longtable}
}

\textbf{HU-24: Consulta de Vigencia de Garantía}

{\def\LTcaptype{none} % do not increment counter
\begin{longtable}[]{@{}
  >{\raggedright\arraybackslash}p{(\linewidth - 2\tabcolsep) * \real{0.2778}}
  >{\raggedright\arraybackslash}p{(\linewidth - 2\tabcolsep) * \real{0.7222}}@{}}
\toprule\noalign{}
\begin{minipage}[b]{\linewidth}\raggedright
\textbf{Campo}
\end{minipage} & \begin{minipage}[b]{\linewidth}\raggedright
\textbf{Contenido}
\end{minipage} \\
\midrule\noalign{}
\endhead
\bottomrule\noalign{}
\endlastfoot
ID e issue & HU-24 \\
Historia & Como cliente, quiero consultar la vigencia de la garantía de
30 días en mano de obra asociada a mi servicio anterior, para reclamar
ajustes en caso de incidencias. \\
Criterio de aceptación 1 & Dado que el cliente revisa el expediente de
su servicio pasado en el portal web, cuando consulta el apartado de
garantías, entonces visualiza la fecha exacta de expiración (30 días
naturales en mano de obra). \\
Criterio de aceptación 2 & Dado que la garantía se encuentra vigente,
cuando se solicita una revisión por el mismo concepto, entonces el
sistema valida el estatus sin generar cargos adicionales de mano de
obra. \\
Reglas de negocio & RN-13 (Garantía de mano de obra). \\
Wireframe & Se define en 4.10 (Review 2) \\
Casos de prueba & Se definen en 5.2 (Review 2) \\
\end{longtable}
}

\textbf{HU-25: Operación Offline del Personal Interno}

{\def\LTcaptype{none} % do not increment counter
\begin{longtable}[]{@{}
  >{\raggedright\arraybackslash}p{(\linewidth - 2\tabcolsep) * \real{0.2778}}
  >{\raggedright\arraybackslash}p{(\linewidth - 2\tabcolsep) * \real{0.7222}}@{}}
\toprule\noalign{}
\begin{minipage}[b]{\linewidth}\raggedright
\textbf{Campo}
\end{minipage} & \begin{minipage}[b]{\linewidth}\raggedright
\textbf{Contenido}
\end{minipage} \\
\midrule\noalign{}
\endhead
\bottomrule\noalign{}
\endlastfoot
ID e issue & HU-25 \\
Historia & Como recepcionista o jefe de mecánicos, quiero registrar la
llegada de autos, tomar fotografías y hacer diagnósticos sin conexión a
internet usando la batería del equipo, para que la operación no se
detenga ante fallas eléctricas. \\
Criterio de aceptación 1 & Dado que ocurre un corte de energía o caída
de internet en el taller, cuando el personal opera mediante las baterías
de laptops o tabletas, entonces el sistema local permite continuar
capturando registros, checklists y fotografías. \\
Criterio de aceptación 2 & Dado que se trabaja de forma local (offline),
cuando se guardan los datos, entonces el almacenamiento temporal
mediante Service Workers e IndexedDB asegura que no exista pérdida de
información. \\
Reglas de negocio & Sin regla de negocio específica; ver RNF-06 y
RNF-08. \\
Wireframe & Se define en 4.10 (Review 2) \\
Casos de prueba & Se definen en 5.2 (Review 2) \\
\end{longtable}
}

\textbf{HU-26: Sincronización Automática en la Nube}

{\def\LTcaptype{none} % do not increment counter
\begin{longtable}[]{@{}
  >{\raggedright\arraybackslash}p{(\linewidth - 2\tabcolsep) * \real{0.2778}}
  >{\raggedright\arraybackslash}p{(\linewidth - 2\tabcolsep) * \real{0.7222}}@{}}
\toprule\noalign{}
\begin{minipage}[b]{\linewidth}\raggedright
\textbf{Campo}
\end{minipage} & \begin{minipage}[b]{\linewidth}\raggedright
\textbf{Contenido}
\end{minipage} \\
\midrule\noalign{}
\endhead
\bottomrule\noalign{}
\endlastfoot
ID e issue & HU-26 \\
Historia & Como recepcionista o jefe de mecánicos, quiero que los
registros y fotografías que capturé sin conexión se sincronicen
automáticamente cuando vuelva la red, para no perder ni duplicar
información. \\
Criterio de aceptación 1 & Dado que se restablece la conexión a internet
en el taller, cuando el sistema detecta conectividad, entonces inicia de
forma automática la sincronización de los datos almacenados
localmente. \\
Criterio de aceptación 2 & Dado que finaliza el proceso de
sincronización, cuando se actualiza la base de datos central en la nube,
entonces la información multimedia y de auditoría se consolida sin
duplicidades. \\
Reglas de negocio & Sin regla de negocio específica; ver RNF-09. \\
Wireframe & Se define en 4.10 (Review 2) \\
Casos de prueba & Se definen en 5.2 (Review 2) \\
\end{longtable}
}

\textbf{HU-27: Gestión de Roles y Permisos de Seguridad}

{\def\LTcaptype{none} % do not increment counter
\begin{longtable}[]{@{}
  >{\raggedright\arraybackslash}p{(\linewidth - 2\tabcolsep) * \real{0.2778}}
  >{\raggedright\arraybackslash}p{(\linewidth - 2\tabcolsep) * \real{0.7222}}@{}}
\toprule\noalign{}
\begin{minipage}[b]{\linewidth}\raggedright
\textbf{Campo}
\end{minipage} & \begin{minipage}[b]{\linewidth}\raggedright
\textbf{Contenido}
\end{minipage} \\
\midrule\noalign{}
\endhead
\bottomrule\noalign{}
\endlastfoot
ID e issue & HU-27 \\
Historia & Como gerente del taller, quiero dar de alta al personal y
asignarle su rol (Gerente, Recepcionista, Jefe de Mecánicos o Encargado
de Almacén), para que cada quien acceda solo a la información que le
corresponde. \\
Criterio de aceptación 1 & Dado que un usuario intenta acceder a la
plataforma, cuando introduce sus credenciales, entonces el sistema
autentica su identidad y restringe los menús según su rol asignado
(Gerente, Recepcionista, Jefe de Mecánicos o Encargado de Almacén). \\
Criterio de aceptación 2 & Dado que un usuario sin privilegios intenta
ingresar a una ruta protegida (como reportes gerenciales), entonces el
sistema deniega el acceso y muestra una pantalla de error por permisos
insuficientes. \\
Reglas de negocio & Sin regla de negocio específica; ver RNF-05 y 4.8
(control de acceso por rol). \\
Wireframe & Se define en 4.10 (Review 2) \\
Casos de prueba & CPA-28, CPA-29 \\
\end{longtable}
}

\textbf{HU-28: Accesibilidad PWA Responsiva para Clientes}

{\def\LTcaptype{none} % do not increment counter
\begin{longtable}[]{@{}
  >{\raggedright\arraybackslash}p{(\linewidth - 2\tabcolsep) * \real{0.2778}}
  >{\raggedright\arraybackslash}p{(\linewidth - 2\tabcolsep) * \real{0.7222}}@{}}
\toprule\noalign{}
\begin{minipage}[b]{\linewidth}\raggedright
\textbf{Campo}
\end{minipage} & \begin{minipage}[b]{\linewidth}\raggedright
\textbf{Contenido}
\end{minipage} \\
\midrule\noalign{}
\endhead
\bottomrule\noalign{}
\endlastfoot
ID e issue & HU-28 \\
Historia & Como cliente, quiero visualizar la plataforma web desde mi
teléfono celular de forma responsiva y sin instalar aplicaciones
adicionales, para consultar mi estatus con facilidad. \\
Criterio de aceptación 1 & Dado que el cliente abre el enlace recibido
por SMS o correo en un navegador móvil (iOS/Android), cuando carga la
página, entonces el diseño se adapta responsivamente a la pantalla sin
requerir descargas de tiendas de aplicaciones. \\
Criterio de aceptación 2 & Dado que el cliente navega por el portal web,
cuando consulta el progreso de su vehículo o galería de fotos, entonces
los elementos táctiles responden de manera fluida y adaptada a
dispositivos móviles. \\
Reglas de negocio & RN-01 (Expiración y vigencia de enlaces).
Portabilidad: RNF-01. \\
Wireframe & W-04 (4.10, Review 2) \\
Casos de prueba & Se definen en 5.2 (Review 2) \\
\end{longtable}
}

\textbf{HU-29: Generación de Reportes Gerenciales}

{\def\LTcaptype{none} % do not increment counter
\begin{longtable}[]{@{}
  >{\raggedright\arraybackslash}p{(\linewidth - 2\tabcolsep) * \real{0.2778}}
  >{\raggedright\arraybackslash}p{(\linewidth - 2\tabcolsep) * \real{0.7222}}@{}}
\toprule\noalign{}
\begin{minipage}[b]{\linewidth}\raggedright
\textbf{Campo}
\end{minipage} & \begin{minipage}[b]{\linewidth}\raggedright
\textbf{Contenido}
\end{minipage} \\
\midrule\noalign{}
\endhead
\bottomrule\noalign{}
\endlastfoot
ID e issue & HU-29 \\
Historia & Como gerente, quiero generar reportes mensuales de ingresos,
servicios realizados y tiempos de ocupación de las 4 bahías, para
evaluar el rendimiento general del taller. \\
Criterio de aceptación 1 & Dado que el gerente ingresa al módulo de
reportes, cuando selecciona el intervalo de fechas mensual, entonces el
sistema genera métricas consolidadas de ingresos y servicios
completados. \\
Criterio de aceptación 2 & Dado que se consultan los tiempos operativos,
cuando el sistema procesa las órdenes, entonces muestra gráficos de
ocupación de las 4 bahías para medir la eficiencia del taller. \\
Reglas de negocio & Sin regla de negocio específica; ver RF-11. \\
Wireframe & Se define en 4.10 (Review 2) \\
Casos de prueba & Se definen en 5.2 (Review 2) \\
\end{longtable}
}

\textbf{HU-30: Cumplimiento de Privacidad y Uso de Datos Ficticios}

{\def\LTcaptype{none} % do not increment counter
\begin{longtable}[]{@{}
  >{\raggedright\arraybackslash}p{(\linewidth - 2\tabcolsep) * \real{0.2778}}
  >{\raggedright\arraybackslash}p{(\linewidth - 2\tabcolsep) * \real{0.7222}}@{}}
\toprule\noalign{}
\begin{minipage}[b]{\linewidth}\raggedright
\textbf{Campo}
\end{minipage} & \begin{minipage}[b]{\linewidth}\raggedright
\textbf{Contenido}
\end{minipage} \\
\midrule\noalign{}
\endhead
\bottomrule\noalign{}
\endlastfoot
ID e issue & HU-30 \\
Historia & Como cliente del taller, quiero conocer el aviso de
privacidad antes de proporcionar mis datos, para saber para qué se
usarán conforme a la LFPDPPP. \\
Criterio de aceptación 1 & Dado que la recepcionista registra a un
cliente nuevo, cuando captura sus datos personales, entonces el sistema
muestra el aviso de privacidad simplificado (4.8) antes de guardar el
registro. \\
Criterio de aceptación 2 & Dado que el cliente abre su enlace seguro por
primera vez, cuando carga la página, entonces el sistema muestra el
enlace al aviso de privacidad integral y cómo ejercer sus derechos
ARCO. \\
Reglas de negocio & Sin regla de negocio específica; ver 4.8 y
RNF-05. \\
Wireframe & Se define en 4.10 (Review 2) \\
Casos de prueba & CPA-31 \\
\end{longtable}
}

\figura{figuras/story-map.png}{width=\linewidth}{4.1}{Story map con el corte del release v0.1.0: actividades del flujo de atención en columnas, historias ordenadas por prioridad y, bajo la línea, las que pasan a la Fase 2. Elaboración propia}{docs/04-requisitos/story-map.py}

\subsection{4.6 Requisitos no
funcionales}\label{requisitos-no-funcionales}

{\def\LTcaptype{none} % do not increment counter
\begin{longtable}[]{@{}
  >{\raggedright\arraybackslash}p{(\linewidth - 8\tabcolsep) * \real{0.0962}}
  >{\raggedright\arraybackslash}p{(\linewidth - 8\tabcolsep) * \real{0.2136}}
  >{\raggedright\arraybackslash}p{(\linewidth - 8\tabcolsep) * \real{0.2735}}
  >{\raggedright\arraybackslash}p{(\linewidth - 8\tabcolsep) * \real{0.2137}}
  >{\raggedright\arraybackslash}p{(\linewidth - 8\tabcolsep) * \real{0.2031}}@{}}
\toprule\noalign{}
\begin{minipage}[b]{\linewidth}\raggedright
\textbf{ID}
\end{minipage} & \begin{minipage}[b]{\linewidth}\raggedright
\textbf{Característica ISO/IEC 25010:2023}
\end{minipage} & \begin{minipage}[b]{\linewidth}\raggedright
\textbf{Requisito}
\end{minipage} & \begin{minipage}[b]{\linewidth}\raggedright
\textbf{Métrica y umbral}
\end{minipage} & \begin{minipage}[b]{\linewidth}\raggedright
\textbf{Cómo se verificará}
\end{minipage} \\
\midrule\noalign{}
\endhead
\bottomrule\noalign{}
\endlastfoot
RNF-01 & Flexibilidad (instalabilidad y adaptabilidad) & El portal del
cliente deberá ser una PWA responsiva que funcione en navegadores
móviles iOS y Android sin instalar software ni conceder permisos
especiales. & Funciona sin errores bloqueantes en las 2 últimas
versiones de Safari (iOS) y Chrome (Android); 0 permisos solicitados al
cliente; sin desplazamiento horizontal en anchos de 360 a 430 px. &
Pruebas manuales en 2 móviles reales (iOS y Android) con una lista de
verificación de las pantallas del cliente y revisión con las
herramientas de desarrollo del navegador. \\
RNF-02 & Capacidad de interacción & La interfaz de la tableta del jefe
de Mecánicos deberá tener botones grandes, alto contraste y controles
táctiles adaptados al uso con manos sucias y en ambiente ruidoso. &
Botones de al menos 48 × 48 px con separación mínima de 8 px; contraste
de texto de al menos 7:1; al menos 95 \% de aciertos en 20 toques con
guantes de trabajo. & Auditoría con Lighthouse o axe DevTools
(gratuitos), medición del contraste con un verificador y prueba en la
tableta real con guantes. \\
RNF-03 & Capacidad de interacción & La interfaz de la tableta deberá
ofrecer flujos de pocos pasos para las operaciones frecuentes del
taller. & Asignar bahía, actualizar el estado de una orden y sellar un
trabajo en máximo 4 toques y 30 s cada una. & Prueba de usabilidad
cronometrada con al menos 3 usuarios (o compañeros que simulen el rol),
repitiendo cada operación 3 veces. \\
RNF-04 & Seguridad (no repudio e integridad) & Toda autorización de
presupuesto o de sobrecosto mayor al 10 \% deberá registrar marca de
tiempo y metadatos inalterables, incluida la dirección IP, para evitar
disputas con el cliente. & 100 \% de las autorizaciones guardan fecha y
hora UTC, IP, identificador de orden y huella SHA-256 del documento; 100
\% de los intentos de alterar un registro se detectan al verificar la
huella. & Prueba con 20 autorizaciones, revisión de los registros
guardados e intento de modificar directamente la base de datos para
confirmar que la verificación lo detecta. \\
RNF-05 & Seguridad (confidencialidad y cumplimiento legal) & El manejo
de datos personales deberá apegarse a la Ley Federal de Protección de
Datos Personales en Posesión de los Particulares y usar solo datos
ficticios en las pruebas. & 100 \% de las conexiones con HTTPS (TLS 1.2
o superior); datos personales sensibles cifrados en reposo; aviso de
privacidad visible en el 100 \% de los formularios que capturan datos; 0
registros reales en el ambiente de pruebas. & Análisis con testssl.sh o
SSL Labs, lista de verificación de la ley revisada por el equipo y
revisión de los datos de prueba antes de cada entrega. \\
RNF-06 & Eficiencia de desempeño & El registro de recepción se deberá
guardar de forma local con rapidez, aun sin conexión. & Guardar un
registro con 5 fotografías en máximo 2 s (percentil 95) en la tableta o
laptop del taller. & 20 repeticiones cronometradas con la herramienta de
rendimiento del navegador, calculando el percentil 95. \\
RNF-07 & Mantenibilidad (modificabilidad) & El desarrollo deberá usar
únicamente herramientas de código abierto, frameworks de uso libre (como
FastAPI y Python) y planes educativos. & 100 \% de las dependencias con
licencia de código abierto o plan educativo vigente; costo de
licenciamiento de \$0 en esta fase. & Inventario de dependencias con
herramientas de revisión de licencias (por ejemplo, pip-licenses) y
registro de los planes educativos usados. \\
RNF-08 & Fiabilidad (tolerancia a fallos) & El personal interno deberá
poder operar sin internet o sin energía eléctrica, con equipos en
batería, bajo una arquitectura local con sincronización diferida. & Las
7 operaciones definidas (registro de entrada, checklist, fotografías,
diagnósticos, asignación de bahías, consulta de proveedores y
autorizaciones gerenciales) funcionan al 100 \% durante al menos 4 h sin
conexión; 0 datos perdidos. & Casos de prueba de las 7 operaciones con
la red desconectada y un equipo con batería, comparando los registros al
terminar. \\
RNF-09 & Fiabilidad (recuperabilidad) & Al restablecerse la conexión, el
sistema deberá sincronizar automáticamente datos, archivos multimedia y
registros de auditoría con la base de datos central, sin duplicados. &
100 \% de los registros sincronizados, 0 duplicados y 0 pérdidas; 200
registros con 50 fotografías en máximo 5 min. & Crear 200 registros sin
conexión, reconectar y comparar conteos y huellas con la base central;
repetir 3 veces e interrumpir la conexión a mitad de una
sincronización. \\
RNF-10 & Fiabilidad (tolerancia a fallos del cliente) & Ante una falla
de conectividad durante la consulta de estatus o la autorización 2FA, la
interfaz del cliente deberá informar el estado y reintentar
automáticamente. & Mensaje de estado en máximo 3 s; reintento automático
cada 30 s o menos hasta completar; 0 autorizaciones duplicadas o
perdidas. & Cortar el servidor durante la consulta y durante el envío
del código, cronometrar el mensaje y revisar la base de datos tras
restablecer el servicio. \\
\end{longtable}
}

\subsection{4.7 Reglas de negocio y requisitos de
dominio}\label{reglas-de-negocio-y-requisitos-de-dominio}

{\def\LTcaptype{none} % do not increment counter
\begin{longtable}[]{@{}
  >{\raggedright\arraybackslash}p{(\linewidth - 6\tabcolsep) * \real{0.0963}}
  >{\raggedright\arraybackslash}p{(\linewidth - 6\tabcolsep) * \real{0.4871}}
  >{\raggedright\arraybackslash}p{(\linewidth - 6\tabcolsep) * \real{0.2351}}
  >{\raggedright\arraybackslash}p{(\linewidth - 6\tabcolsep) * \real{0.1815}}@{}}
\toprule\noalign{}
\begin{minipage}[b]{\linewidth}\raggedright
\textbf{ID}
\end{minipage} & \begin{minipage}[b]{\linewidth}\raggedright
\textbf{Regla de negocio}
\end{minipage} & \begin{minipage}[b]{\linewidth}\raggedright
\textbf{Fuente (anteproyecto / entrevista con fecha)}
\end{minipage} & \begin{minipage}[b]{\linewidth}\raggedright
\textbf{Requisitos que afecta}
\end{minipage} \\
\midrule\noalign{}
\endhead
\bottomrule\noalign{}
\endlastfoot
RN-01 & Expiración y vigencia de enlaces: El enlace web de seguimiento
permanecerá activo durante la estancia del vehículo en el taller.
Expirará automáticamente al entregar y saldar la unidad. Si el cliente
requiere ver su historial posteriormente, deberá solicitar la generación
de un nuevo enlace. La vigencia del enlace es independiente de la
vigencia de la cotización (RN-11). & Entrevista con el cliente (19/10/2026), minuta A.2, E-03 & RF-02.1, RF-02.2 \\
RN-02 & Doble verificación en autorizaciones: Toda aprobación remota o
presencial de cotización exigirá una doble verificación (código 2FA de 6
dígitos, de un solo uso y vigente 5 minutos, enviado al correo
electrónico registrado del cliente) para evitar reclamos por trabajos no
autorizados. & Entrevista con el cliente (19/10/2026), minuta A.2, E-01 & RF-06, RNF-04 \\
RN-03 & Selección granular de rubros: El cliente puede aprobar o
rechazar de manera independiente los conceptos opcionales desglosados en
el presupuesto (mano de obra o refacciones), ejecutando únicamente lo
autorizado explícitamente. & Entrevista con el cliente (19/10/2026), minuta A.2, E-02 &
RF-06 \\
RN-04 & Re-cotización obligatoria por sobrecostos (\textgreater10\%):
Cualquier trabajo o repuesto imprevisto durante la reparación que
incremente el costo inicial más de un 10\% detendrá el avance hasta
contar con una nueva autorización validada por el cliente. &
Anteproyecto y propuesta técnica & RF-07, RNF-04 \\
RN-05 & Capacidad física del taller\textbf{:} La atención simultánea
está limitada estrictamente a 4 bahías físicas (1 vehículo por bahía) y
un área de estacionamiento de recepción/espera con capacidad máxima para
5 vehículos. & Anteproyecto y entrevista con el cliente (19/10/2026), minuta A.2, E-04 & RF-05,
RF-09 \\
RN-06 & Reubicación por espera de refacciones: Si un vehículo requiere
refacciones externas con demora de entrega de varios días, se liberará
la bahía de trabajo reubicando la unidad en el estacionamiento hasta el
arribo de las piezas. & Entrevista con el cliente (19/10/2026), minuta A.2, E-05 & RF-05,
RF-08.2 \\
RN-07 & Clasificación de vehículos autorizados: El sistema operará
exclusivamente con vehículos automotrices ligeros (automóviles
particulares, motocicletas y camionetas hasta de 3 toneladas). Se
excluyen bicicletas, tráilers y maquinaria pesada. & Entrevista con el cliente (19/10/2026), minuta A.2, E-06 & RF-03.1 \\
RN-08 & Confirmación manual de citas web: Las solicitudes de cita
agendadas desde el portal público ingresarán en estado "Pendiente" y no
reservarán espacio en la agenda hasta que la recepcionista las confirme
manualmente. & Entrevista con el cliente (19/10/2026), minuta A.2, E-07 & RF-09 \\
RN-09 & Asignación y colaboración de mecánicos\textbf{:} Cada orden
mantendrá un mecánico principal responsable. Otros mecánicos podrán
colaborar en el trabajo técnico sin alterar la titularidad de la orden,
salvo reasignación formal. & Entrevista con el cliente (19/10/2026), minuta A.2, E-08 &
RF-05 \\
RN-10 & Restricción de entrega y excepción a flotillas: Queda prohibida
la entrega física del vehículo con saldo pendiente. Como única
excepción, el Gerente podrá autorizar la entrega a crédito para
vehículos pertenecientes a flotillas empresariales registradas. &
Anteproyecto y entrevista con el cliente (19/10/2026), minuta A.2, E-09 & RF-10 \\
RN-11 & Vigencia de cotizaciones: Toda cotización tiene vigencia de 5
días naturales desde su emisión. Una cotización vencida no puede
autorizarse; el taller debe emitir una nueva versión. & Anteproyecto &
RF-06 \\
RN-12 & Evidencia fotográfica de recepción: Al recibir un vehículo se
exigen el checklist completo y al menos 4 fotografías (frente, trasera,
costado izquierdo y costado derecho); sin ellas la orden no puede pasar
a diagnóstico. & Anteproyecto; el mínimo de 4 fotografías debe validarse
con el cliente & RF-04 \\
RN-13 & Garantía de mano de obra: Los trabajos tienen garantía de 30
días naturales sobre la mano de obra, contados desde la entrega del
vehículo. & RF-03.2; confirmar la fuente con el cliente & RF-03.2 \\
RN-14 & Medidor de combustible descompuesto: Si el indicador de combustible
del tablero no funciona, el nivel no se deja en blanco ni se adivina: se
registra como «No verificable», junto con el nivel que declara el cliente,
el método de estimación usado (4.7.1) y la aceptación del cliente. El
taller no responde por diferencias de combustible en una recepción
marcada como «No verificable», salvo lo que él mismo haya medido. &
Propuesta del equipo; validar con el cliente & RF-04,
HU-01 \\
RN-15 & Operación de contingencia en papel: Si la falta de luz o de
internet dura más que la autonomía de los equipos (RNF-08), la operación
continúa en la Bitácora de Contingencia impresa (4.7.1) con folios
provisionales. Todo lo registrado en papel se captura en el sistema a más
tardar 24 horas después de restablecerse el servicio y conserva la hoja
física escaneada como evidencia. & Propuesta del equipo; amplía HU-25;
validar con el cliente & RNF-08, RNF-09, RF-04, RF-06,
RF-10 \\
\end{longtable}
}

\subsubsection{4.7.1 Procedimientos de
contingencia}\label{procedimientos-de-contingencia}

Estos procedimientos detallan RN-14 y RN-15. Son operativos: indican qué
hace el personal cuando falla el vehículo del cliente o la infraestructura
del taller, y qué tendrá que soportar el sistema para que esos casos queden
registrados.

\paragraph{A. Vehículo con gasolina pero con el medidor descompuesto
(RN-14).}

El nivel de combustible al ingreso protege a ambas partes ante reclamos
(«dejé medio tanque y me lo entregaron en reserva»). Cuando el indicador
del tablero no funciona, la recepcionista aplica, en este orden, el primer
método que esté disponible y anota cuál usó:

\begin{enumerate}
\tightlist
\item \textbf{Lectura por OBD-II.} Muchas veces solo falla la aguja o el
tablero, pero el sensor (flotador) del tanque sigue funcionando. Con el
escáner OBD-II del taller se lee el parámetro «Fuel Tank Level Input» (PID
\texttt{0x2F}), que da el porcentaje del tanque, y se convierte al selector
de 5 niveles (0--12\,\% Reserva, 13--37\,\% 1/4, 38--62\,\% 1/2,
63--87\,\% 3/4, 88--100\,\% Lleno). Se fotografía la pantalla del escáner.
\item \textbf{Computadora de viaje.} Si el vehículo muestra la autonomía
restante (km), se fotografía y se registra como referencia, aclarando que
es una estimación del propio vehículo.
\item \textbf{Medición física, solo si el trabajo lo requiere.} Si la
reparación implica bajar o vaciar el tanque (bomba de gasolina, sensor de
nivel, fugas), el combustible se extrae con la bomba de trasvase a bidones
graduados aprobados para combustible y se registran los \textbf{litros}
medidos, con fotografía del bidón. Al terminar, ese mismo volumen se
regresa al tanque. No se abre el sistema de combustible solo para medir:
es un riesgo de incendio y no forma parte del servicio autorizado.
\item \textbf{Declaración del cliente.} Si nada de lo anterior es posible,
se registra el nivel que dice el cliente con la leyenda «Declarado por el
cliente, no verificable».
\end{enumerate}

En todos los casos se fotografía el tablero encendido (evidencia de que el
indicador no funciona), se agrega al checklist el hallazgo «Indicador de
combustible inoperante» y se le ofrece al cliente diagnosticarlo como
servicio adicional. El cliente acepta lo registrado con su firma en la
recepción. Al entregar, el nivel se vuelve a medir con el mismo método para
que la comparación sea justa. La conversión del porcentaje OBD-II a los 5
niveles es una de las decisiones candidatas a modelarse en DMN (7.9).

\emph{Impacto en el sistema:} el selector de combustible del wireframe W-01
gana la opción «No verificable», que obliga a capturar el método de
estimación (OBD-II, computadora de viaje, medición en litros o declaración
del cliente), el valor obtenido y una fotografía del tablero.
\pendiente{agregar los campos \texttt{metodo\_combustible} y
\texttt{litros\_medidos} a RecepcionVehicular (6.2) y su caso de prueba en
la sección 5.}

\paragraph{B. Cortes prolongados de luz o de internet (RN-15).}

La operación sin conexión de la PWA (HU-25, RNF-08) cubre los cortes
cortos. Para los largos se definen tres niveles:

{\small
\begin{longtable}{@{}L{0.17\linewidth}L{0.33\linewidth}L{0.44\linewidth}@{}}
\toprule
\textbf{Nivel} & \textbf{Situación} & \textbf{Cómo se opera} \\
\midrule
\endhead
\bottomrule
\endlastfoot
1. Sin internet & Hay luz, pero no hay conexión. & La PWA trabaja con
IndexedDB y sincroniza al volver la conexión (RNF-09). El cliente no
recibe enlace ni puede autorizar en línea; se le avisa por teléfono. \\
2. Sin luz & Los equipos funcionan con batería. & Igual que el nivel 1. Se
apagan los equipos que no se usen y se reservan las baterías para
recepción y diagnóstico. \\
3. Contingencia en papel & El corte supera 4 horas o la batería de los
equipos baja del 20\,\%. & El Gerente declara la contingencia y anota la
hora de inicio en la Bitácora. Todo se registra en papel, como se describe
abajo. \\
\end{longtable}
}

La \textbf{Bitácora de Contingencia} es una libreta o carpeta con formatos
impresos y numerados de antemano, que se guarda en recepción. Contiene:

\begin{itemize}
\tightlist
\item \textbf{Hoja de recepción}: folio provisional \texttt{C-AAAAMMDD-NN}
(consecutivo del día, escrito a mano), datos del cliente y del vehículo,
kilometraje, nivel de combustible (con la opción «No verificable» de la
parte A), checklist de carrocería e interiores y firma del cliente. Las 4
fotografías obligatorias (RN-12) se toman con el celular del taller y se
nombran con el folio provisional.
\item \textbf{Hoja de diagnóstico y cotización}, con el desglose de mano de
obra y refacciones.
\item \textbf{Hoja de autorización}: como no es posible enviar el código
2FA por correo (RN-02), la única autorización válida en contingencia es la
\textbf{firma autógrafa del cliente en presencia del personal}, con fecha y
hora. No se aceptan autorizaciones por teléfono.
\item \textbf{Recibo provisional} numerado para los pagos (solo efectivo,
porque la terminal bancaria tampoco funciona). Sigue vigente RN-10: no se
entrega un vehículo con saldo pendiente.
\item \textbf{Registro de movimientos}: cambios de estado y de bahía, con
hora y nombre de quien los hizo.
\end{itemize}

\textbf{Al restablecerse el servicio}, el Gerente anota la hora de fin en
la Bitácora. En las 24 horas siguientes, la recepcionista captura cada hoja
en el sistema: la orden recibe su folio definitivo OS-AAAA-NNNN y conserva
el folio provisional como referencia, se adjuntan las fotografías y la hoja
escaneada, y las autorizaciones quedan registradas con medio «Firma en
papel». El Gerente revisa que no falte ningún folio provisional del
consecutivo y firma la Bitácora como conciliada. Las hojas físicas se
resguardan con el mismo cuidado que los datos personales del sistema (4.8).

\pendiente{agregar al modelo el medio de autorización «Firma en papel» y el
campo \texttt{folio\_provisional}; diseñar e imprimir los formatos de la
Bitácora; validar con el cliente el umbral de 4 horas y el plazo de 24
horas.}

